\chapter{Simulacro C: compresión axial y radial}
\noindent\textbf{Condiciones.} 120 minutos orientativos. Declare convenios, hipótesis y unidades; el sentido positivo de las velocidades tangenciales es el de la velocidad periférica $U$. Gas perfecto con $c_p=1005\,\mathrm{J/(kg\,K)}$ y $\gamma=1.4$. Propuesta original, sin carácter oficial.
\section{Cuestiones razonadas (2 puntos)}
Explique por qué el trabajo de Euler de un escalón no depende de la eficiencia y sí lo hace la razón de presiones. Indique qué mide el cociente de velocidades relativas $W_2/W_1$ y por qué se limita en un rotor de compresor. Explique por qué el deslizamiento reduce el trabajo de un rodete radial.
\section{Escalón de compresor axial (4 puntos)}
Un escalón trabaja al radio medio con $T_{01}=288.15\,\mathrm K$, $p_{01}=101325\,\mathrm{Pa}$, $U=250\,\mathrm{m/s}$, $C_x=160\,\mathrm{m/s}$ constante, remolino absoluto $C_{\theta1}=70\,\mathrm{m/s}$ a la entrada del rotor y $C_{\theta2}=180\,\mathrm{m/s}$ a la salida. El rendimiento total-total es $0.88$ y el caudal $20\,\mathrm{kg/s}$.
\begin{enumerate}[label=\alph*)]
\item Calcule el trabajo específico, el aumento de temperatura total, la razón de presiones y la presión total de salida. (1.5)
\item Calcule la potencia absorbida y la reacción del escalón. (1.0)
\item Calcule las velocidades relativas, los ángulos relativos y el cociente de velocidades relativas; comente su valor. (1.5)
\end{enumerate}
\section{Rodete centrífugo (4 puntos)}
Un rodete de $0.30\,\mathrm m$ de diámetro gira a $36000\,\mathrm{min^{-1}}$ sin prerrotación y con deslizamiento $0.85$. La temperatura total de entrada es $300\,\mathrm K$ y el rendimiento total-total $0.80$.
\begin{enumerate}[label=\alph*)]
\item Calcule la velocidad periférica de salida y el remolino absoluto de salida. (1.0)
\item Calcule el trabajo específico, el aumento de temperatura total y la razón de presiones. (1.5)
\item Calcule el Mach periférico y la potencia para $1.2\,\mathrm{kg/s}$; discuta qué límites de diseño aparecen. (1.5)
\end{enumerate}

\chapter{Simulacro D: turbina axial, pala y refrigeración}
\noindent\textbf{Condiciones.} 120 minutos orientativos. Declare convenios, hipótesis y unidades. Gas caliente con $c_p=1150\,\mathrm{J/(kg\,K)}$ y $\gamma=1.33$. Propuesta original, sin carácter oficial.
\section{Cuestiones razonadas (2 puntos)}
Defina coeficiente de carga y de flujo y explique cómo condicionan el diseño del escalón. Explique por qué la reacción de una turbina se define con la caída de entalpía estática del rotor y qué hipótesis hace falta para igualar el trabajo con la caída del escalón. Indique por qué la tensión centrífuga de una pala no depende de la forma de la sección en el modelo de barra uniforme.
\section{Escalón de turbina con descarga axial (4 puntos)}
Un escalón repetitivo trabaja al radio medio con $T_{01}=1400\,\mathrm K$, $U=320\,\mathrm{m/s}$ y $C_x=200\,\mathrm{m/s}$ constante. El remolino a la salida del estator es $C_{\theta2}=384\,\mathrm{m/s}$ y el escalón descarga axialmente. El rendimiento total-total es $0.90$.
\begin{enumerate}[label=\alph*)]
\item Calcule el trabajo específico, los coeficientes de carga y de flujo y la caída de temperatura total. (1.5)
\item Calcule la razón de expansión total-total y el ángulo absoluto de entrada al rotor. (1.0)
\item Calcule las velocidades relativas al cuadrado a la entrada y a la salida del rotor y la reacción. (1.5)
\end{enumerate}
\section{Pala y pared refrigerada (4 puntos)}
Una pala de acero de densidad $8200\,\mathrm{kg/m^3}$ gira a $9000\,\mathrm{min^{-1}}$ con cubo a $0.35\,\mathrm m$ y punta a $0.45\,\mathrm m$. Se adopta una tensión admisible docente de $400\,\mathrm{MPa}$. La pared de la pala tiene $1.5\,\mathrm{mm}$ de espesor y conductividad $25\,\mathrm{W/(m\,K)}$; el gas exterior ($T_{aw}=1500\,\mathrm K$, $h_g=1800\,\mathrm{W/(m^2\,K)}$) y el refrigerante ($700\,\mathrm K$, $h_c=2500\,\mathrm{W/(m^2\,K)}$) dan una red estacionaria unidimensional.
\begin{enumerate}[label=\alph*)]
\item Calcule la tensión centrífuga de barra uniforme, la velocidad en la punta y el coeficiente de seguridad. (1.5)
\item Calcule el flujo de calor, la temperatura de la pared caliente y la diferencia de temperatura a través de la pared. (1.5)
\item Calcule la tensión térmica con restricción total uniaxial, con $E=190\,\mathrm{GPa}$ y $\alpha=14\times10^{-6}\,\mathrm{K^{-1}}$, e indique qué fenómenos deja fuera. (1.0)
\end{enumerate}

\chapter{Simulacro E: fuera de diseño y acoplamiento}
\noindent\textbf{Condiciones.} 120 minutos orientativos. Declare convenios, hipótesis y unidades. Temperatura y presión de referencia $288.15\,\mathrm K$ y $101325\,\mathrm{Pa}$. Propuesta original, sin carácter oficial.
\section{Cuestiones razonadas (2 puntos)}
Explique qué representa el régimen corregido y bajo qué hipótesis dos puntos con el mismo régimen y caudal corregidos tienen la misma razón de presiones. Indique por qué un margen de bombeo no puede deducirse solo de los datos de un punto de funcionamiento. Explique por qué la potencia de la turbina de un generador de gas debe igualar la del compresor más las pérdidas mecánicas.
\section{Corrección de un punto de ensayo (4 puntos)}
Un compresor se ensaya con $T_{in}=250\,\mathrm K$ y $p_{in}=70000\,\mathrm{Pa}$. Mide un caudal de $14\,\mathrm{kg/s}$ a $10200\,\mathrm{min^{-1}}$.
\begin{enumerate}[label=\alph*)]
\item Calcule el régimen y el caudal corregidos. (2.0)
\item Determine el régimen físico y el caudal físico del punto equivalente con $T=320\,\mathrm K$ y $p=101325\,\mathrm{Pa}$. (1.5)
\item Indique qué hipótesis de semejanza se usan. (0.5)
\end{enumerate}
\section{Generador de gas (4 puntos)}
El compresor de un generador de gas mueve $22\,\mathrm{kg/s}$ con un aumento de temperatura total de $240\,\mathrm K$ ($c_p=1005\,\mathrm{J/(kg\,K)}$). La turbina recibe $(1+0.02-0.05)$ veces ese caudal (combustible $2\,\%$, sangrado $5\,\%$), con $c_p=1150\,\mathrm{J/(kg\,K)}$, $\gamma=1.33$, rendimiento mecánico $0.99$, rendimiento total-total $0.88$ y temperatura total de entrada $1450\,\mathrm K$.
\begin{enumerate}[label=\alph*)]
\item Calcule la potencia del compresor. (1.0)
\item Calcule la caída de temperatura total de la turbina. (1.5)
\item Calcule la razón de expansión total-total e indique tres efectos que el modelo omite. (1.5)
\end{enumerate}
